\chapter*{도서 기획안 및 상세 목차}
\markboth{도서 기획안 및 상세 목차}{도서 기획안 및 상세 목차}
\addcontentsline{toc}{chapter}{도서 기획안 및 상세 목차}

\section*{1. 도서 기획 개요}

\subsection*{1.1 도서 콘셉트 및 차별화 전략}
\textbf{본 저서의 핵심 콘셉트}는 ``이론과 실무의 가교(Bridge between Theory and Practice)''입니다. 기존 AI 거버넌스 관련 서적들이 대체로 추상적인 원칙 나열이나 해외 사례 번역에 그치는 반면, 본 저서는 데이터 패브릭 아키텍처, 공공부문 및 헬스케어 도메인 프로젝트 경험을 바탕으로 한국 기업 환경에 최적화된 실행 가능한 프레임워크를 제시합니다.

\begin{itemize}
 \item \textbf{차별화 요소 1: 계층별 접근법(Tiered Approach)}\\
 경영진용 전략적 관점, 중간관리자용 운영적 관점, 실무자용 기술적 관점을 명확히 분리하여 독자가 자신의 역할을 맞는 내용을 선택적으로 활용할 수 있도록 구성합니다.
 \item \textbf{차별화 요소 2: 컨설팅 방법론 내재화}\\
 각 장마다 `현황 진단 $\rightarrow$ 갭 분석 $\rightarrow$ 개선 로드맵 $\rightarrow$ 실행 체크리스트'의 구조를 일관되게 적용하여, 컨설턴트가 고객사에 그대로 적용할 수 있는 워크플로우를 제공합니다.
 \item \textbf{차별화 요소 3: 산업별 맞춤 모듈}\\
 공공부문, 헬스케어, 금융, 제조업 등 산업별 규제 환경과 리스크 프로파일이 상이하므로, 각 산업에 특화된 거버넌스 요구사항과 구현 사례를 별도 섹션으로 다룹니다.
\end{itemize}

\subsection*{1.2 목표 독자층}
\begin{center}
\begin{tabular}{|p{0.25\textwidth}|p{0.3\textwidth}|p{0.35\textwidth}|}
\hline
\textbf{독자 유형} & \textbf{활용 목적} & \textbf{핵심 관심사} \\
\hline
C-레벨 경영진 & AI 투자 의사결정, 리스크 관리 & ROI, 규제 리스크, 평판 리스크 \\
\hline
AI/데이터 조직 리더 & 거버넌스 체계 수립 및 운영 & 조직 설계, KPI, 프로세스 표준화 \\
\hline
컨설턴트 & 고객사 AI 거버넌스 구축 지원 & 진단 도구, 방법론, 산업별 베스트 프랙티스 \\
\hline
데이터 사이언티스트/ \newline ML 엔지니어 & 기술적 준수 구현 & 모델 문서화, 설명가능성, MLOps 통합 \\
\hline
법무/컴플라이언스 담당자 & 규제 준수 체계 마련 & EU AI Act, 국내 법규, 감사 대응 \\
\hline
\end{tabular}
\end{center}

\subsection*{1.3 시장 포지셔닝}
현재 국내 AI 거버넌스 관련 서적 시장을 분석하면, 대부분이 (1) 해외 프레임워크 소개 중심이거나, (2) 윤리적 담론에 치중하거나, (3) 특정 기술(예: MLOps)에만 집중하는 경향이 있습니다. 본 저서는 ``전사적 AI 거버넌스 구축을 위한 종합 실무 지침서''라는 포지션을 확보하여, 조직이 AI 거버넌스를 처음 도입하거나 고도화할 때 참조할 수 있는 원스톱 레퍼런스를 지향합니다.

\vspace{0.8cm}
\section*{2. 전체 구성 체계}
본 저서는 5부 20장 구성을 제안하며, 각 부는 AI 거버넌스의 핵심 영역을 체계적으로 다룹니다.
\begin{itemize}
 \item \textbf{[제1부] 기반 구축}: AI 거버넌스의 이론적 토대
 \item \textbf{[제2부] 전략 설계}: 조직 맞춤형 거버넌스 프레임워크
 \item \textbf{[제3부] 운영 체계}: 생명주기 전반의 거버넌스 실행
 \item \textbf{[제4부] 기술 구현}: 거버넌스를 위한 엔지니어링
 \item \textbf{[제5부] 고도화}: 산업별 적용과 미래 전망
 \end{itemize}

\vspace{0.8cm}
\section*{3. 상세 목차}

\subsection*{제1부: 기반 구축 --- AI 거버넌스의 이론적 토대}

\textbf{제1장. AI 거버넌스의 개념과 필요성}
\begin{itemize}
 \item 1.1 AI 거버넌스의 정의와 범위 (Data/IT 거버넌스와의 관계, 구성 요소)
 \item 1.2 AI 거버넌스가 필요한 이유 (비즈니스/기술/규제 관점)
 \item 1.3 AI 거버넌스 부재의 결과 (실패 사례 분석, 손실 정량화)
 \item 1.4 AI 거버넌스 성숙도 모델 (5단계 프레임워크, 자가 진단)
 \item {[실무 도구 1-1] AI 거버넌스 성숙도 자가진단 설문지}
\end{itemize}

\medskip
\textbf{제2장. 글로벌 AI 거버넌스 프레임워크 비교 분석}
\begin{itemize}
 \item 2.1 국제기구 및 표준화 기관의 프레임워크 (OECD, ISO 42001, NIST AI RMF 등)
 \item 2.2 주요국 규제 동향 (EU AI Act, 미국 행정명령, 중국/일본/싱가포르 동향)
 \item 2.3 국내 AI 규제 환경 (AI 기본법, 개인정보보호법 등)
 \item 2.4 규제 수렴과 상호운용성
 \item {[실무 도구 2-1] 글로벌 AI 규제 비교 매트릭스} / {[실무 도구 2-2] EU AI Act 대응 체크리스트}
\end{itemize}

\medskip
\textbf{제3장. AI 윤리 원칙의 운용화(Operationalization)}
\begin{itemize}
 \item 3.1 추상적 원칙에서 측정 가능한 지표로 (공정성, 투명성, 책임성 등)
 \item 3.2 원칙 간 충돌과 트레이드오프 관리 (공정성 vs 정확도 등)
 \item 3.3 윤리 원칙의 문서화와 커뮤니케이션
 \item {[실무 도구 3-1] AI 윤리 원칙 운용화 워크시트} / {[실무 도구 3-2] 원칙 충돌 해결 의사결정 트리}
\end{itemize}

\vspace{0.5cm}
\subsection*{제2부: 전략 설계 --- 조직 맞춤형 거버넌스 프레임워크}

\textbf{제4장. AI 거버넌스 조직 설계}
\begin{itemize}
 \item 4.1 거버넌스 조직 모델의 유형 (중앙집중형/연방형/분산형)
 \item 4.2 핵심 역할과 책임(R\&R) 정의 (CAIO, 윤리위원회 등)
 \item 4.3 거버넌스 위원회 체계
 \item 4.4 RACI 매트릭스 설계
 \item {[실무 도구 4-1] 조직 설계 워크숍 가이드} / {[실무 도구 4-2] 직무 기술서 템플릿}
\end{itemize}

\medskip
\textbf{제5장. AI 정책 체계 수립}
\begin{itemize}
 \item 5.1 정책 아키텍처 설계 (3계층 구조)
 \item 5.2 핵심 정책 문서 개발 (AI 활용 정책, 개발 표준 등)
 \item 5.3 정책 생명주기 관리
 \item {[실무 도구 5-1] AI 정책 체계 템플릿 패키지} / {[실무 도구 5-2] 정책 준수 자가 체크리스트}
\end{itemize}

\medskip
\textbf{제6장. AI 리스크 관리 프레임워크}
\begin{itemize}
 \item 6.1 AI 리스크 분류 체계 (기술/운영/컴플라이언스/전략/윤리)
 \item 6.2 리스크 평가 방법론 (정량/정성 평가, AI 영향 평가)
 \item 6.3 리스크 완화 전략
 \item 6.4 리스크 모니터링과 보고 (KRI, 대시보드)
 \item {[실무 도구 6-1] 리스크 등급 분류 트리} / {[실무 도구 6-2] AIA 템플릿} / {[실무 도구 6-3] 리스크 레지스터}
\end{itemize}

\medskip
\textbf{제7장. AI 거버넌스 로드맵 수립}
\begin{itemize}
 \item 7.1 현항 진단(As-Is Analysis)
 \item 7.2 목표 상태 정의(To-Be 설계)
 \item 7.3 실행 로드맵 개발
 \item 7.4 변화 관리와 커뮤니케이션
 \item {[실무 도구 7-1] 로드맵 워크숍 가이드} / {[실무 도구 7-2] 우선순위 결정 매트릭스}
\end{itemize}

\vspace{0.5cm}
\subsection*{제3부: 운영 체계 --- 생명주기 전반의 거버넌스 실행}

\textbf{제8장. AI 유스케이스 발굴과 사전 승인}
\begin{itemize}
 \item 8.1 AI 적용 적합성 평가
 \item 8.2 유스케이스 제안 및 검토 프로세스
 \item 8.3 AI 프로젝트 포트폴리오 관리
 \item {[실무 도구 8-1] 적합성 평가 체크리스트} / {[실무 도구 8-2] 제안서 템플릿}
\end{itemize}

\medskip
\textbf{제9장. 데이터 거버넌스와 AI 거버넌스의 통합}
\begin{itemize}
 \item 9.1 AI를 위한 데이터 요구사항 (품질, 편향 등)
 \item 9.2 데이터 계보(Lineage)와 출처
 \item 9.3 데이터 접근 제어와 프라이버시
 \item 9.4 외부 데이터 및 사전 학습 모델의 거버넌스
 \item {[실무 도구 9-1] 학습 데이터 품질 평가 체크리스트} / {[실무 도구 9-2] 데이터 편향 탐지 프로토콜}
\end{itemize}

\medskip
\textbf{제10장. 모델 개발 단계의 거버넌스}
\begin{itemize}
 \item 10.1 모델 개발 표준과 문서화 (모델 카드 등)
 \item 10.2 모델 검증과 테스트 (성능/공정성/강건성/보안)
 \item 10.3 설명가능성(XAI) 구현
 \item 10.4 모델 승인과 릴리스 관리
 \item {[실무 도구 10-1] 모델 카드 템플릿} / {[실무 도구 10-2] 검증 체크리스트} / {[실무 도구 10-3] 배포 승인 체크리스트}
\end{itemize}

\medskip
\textbf{제11장. AI 거버넌스 감사 및 인증}
\begin{itemize}
 \item 11.1 AI 거버넌스 감사 개요 (목적, 유형, 범위)
 \item 11.2 내부 감사 수행 (계획, 방법론, 체크리스트)
 \item 11.3 외부 감사 및 인증 (ISO 42001, 알고리즘 감사)
 \item 11.4 감사 대응 및 후속 관리
 \item {[실무 도구 11-1] 감사 계획서 템플릿} / {[실무 도구 11-2] 알고리즘 감사 보고서 템플릿}
\end{itemize}

\medskip
\textbf{제12장. AI 시스템 운영 및 생명주기 관리}
\begin{itemize}
 \item 12.1 MLOps와 거버넌스의 통합
 \item 12.2 모델 모니터링과 드리프트 관리
 \item 12.3 모델 폐기(Retire)와 아카이빙
 \item {[실무 도구 12-1] AI 운영 지표 대시보드} / {[실무 도구 12-2] 모델 폐기 체크리스트}
\end{itemize}

\vspace{0.5cm}
\subsection*{제4부: 기술 구현 --- 거버넌스를 위한 엔지니어링}

\textbf{제13장. AI 시스템 아키텍처와 거버넌스}
\begin{itemize}
 \item 13.1 거버넌스 친화적 아키텍처 설계 원칙
 \item 13.2 RAG 아키텍처의 거버넌스
 \item 13.3 LLM 기반 시스템의 거버넌스 (프롬프트, 가드레일 등)
 \item 13.4 에이전트(Agent) 시스템의 거버넌스
 \item {[실무 도구 13-1] RAG 거버넌스 체크리스트} / {[실무 도구 13-2] LLM 가드레일 템플릿}
\end{itemize}

\medskip
\textbf{제14장. 거버넌스 도구와 플랫폼}
\begin{itemize}
 \item 14.1 모델 레지스트리와 카탈로그
 \item 14.2 실험 추적과 재현성 도구
 \item 14.3 모니터링 및 관측성(Observability) 플랫폼
 \item 14.4 컴플라이언스 자동화 도구
 \item {[실무 도구 14-1] 도구 선정 기준표} / {[실무 도구 14-2] 플랫폼 도입 RFP 템플릿}
\end{itemize}

\medskip
\textbf{제15장. 보안과 프라이버시의 기술적 구현}
\begin{itemize}
 \item 15.1 AI 시스템 보안 위협과 대응 (적대적 공격, 데이터 오염)
 \item 15.2 프라이버시 보호 기술 (차분 프라이버시, 연합 학습 등)
 \item 15.3 접근 제어와 권한 관리
 \item {[실무 도구 15-1] 보안 위협 모델링 워크시트} / {[실무 도구 15-2] 프라이버시 기법 가이드}
\end{itemize}

\vspace{0.5cm}
\subsection*{제5부: 고도화 --- 산업별 적용과 미래 전망}

\textbf{제16장. 헬스케어 AI 거버넌스}
\begin{itemize}
 \item 16.1 의료 AI의 규제 환경 (SaMD, FDA 등)
 \item 16.2 의료 AI 특화 리스크 관리
 \item 16.3 헬스케어 데이터 거버넌스
 \item 16.4 구현 사례 연구
 \item {[실무 도구 16-1] 인허가 로드맵} / {[실무 도구 16-2] 임상 검증 프로토콜}
\end{itemize}

\medskip
\textbf{제17장. 금융 AI 거버넌스}
\begin{itemize}
 \item 17.1 금융권 AI 규제 프레임워크
 \item 17.2 금융 AI 특화 리스크 관리 (공정 대출, 시장 조작 등)
 \item 17.3 금융 AI 감사와 검사 대응
 \item 17.4 구현 사례 연구
 \item {[실무 도구 17-1] 모델 검증 체크리스트} / {[실무 도구 17-2] 공정성 검증 보고서}
\end{itemize}

\medskip
\textbf{제18장. 공공부문 AI 거버넌스}
\begin{itemize}
 \item 18.1 공공 AI의 특수성 (Public AI)
 \item 18.2 공공 AI 규제와 정책
 \item 18.3 공공 AI 조례와 검증
 \item 18.4 구현 사례 연구
 \item {[실무 도구 18-1] 도입 타당성 평가 프레임워크} / {[실무 도구 18-2] 시민 영향 평가 가이드}
\end{itemize}

\medskip
\textbf{제19장. 제조/산업 AI 거버넌스}
\begin{itemize}
 \item 19.1 산업 AI의 특수성 (Safety-Critical, OT/IT 융합)
 \item 19.2 산업 AI 표준과 인증
 \item 19.3 구현 사례 연구
 \item {[실무 도구 19-1] 안전성 평가 체크리스트} / {[실무 도구 19-2] OT-IT 통합 프레임워크}
\end{itemize}

\medskip
\textbf{제20장. AI 거버넌스의 미래와 지속적 진화}
\begin{itemize}
 \item 20.1 신기술과 거버넌스 도전 (AGI, 에이전트, 뉴로심볼릭)
 \item 20.2 규제 환경의 진화 전망
 \item 20.3 거버넌스 조직의 진화 (ESG 통합 등)
 \item 20.4 지속적 개선 체계
 \item {[실무 도구 20-1] KPI 대시보드 가이드} / {[실무 도구 20-2] 연간 검토 체크리스트}
\end{itemize}

\vspace{1cm}
\section*{4. 부록 구성}

\begin{itemize}
 \item \textbf{부록 A.} 용어 사전 (국문-영문 대조)
 \item \textbf{부록 B.} 참조 법규 및 표준 목록 (국내외)
 \item \textbf{부록 C.} 실무 도구 모음 (전자 파일 다운로드 안내)
 \item \textbf{부록 D.} AI 거버넌스 성숙도 평가 도구 상세 버전
 \item \textbf{부록 E.} 추천 학습 자료 및 커뮤니티
\end{itemize}
